Reťazce DNA alebo RNA môžu byť čítané rôznou rýchlosťou počas nanopórového sekvenovania, čo môže ovplyvniť kvalitu a presnosť zaznamenaného signálu \cite{Zheng2023Nanopore}.
Rýchlosť prechodu molekuly cez nanopór je ovplyvnená viacerými faktormi, vrátane veľkosti molekuly, jej konformácie, interakcií s nanopórom a podmienok prostredia (napríklad teplota a iónová sila) \cite{Zheng2023Nanopore}.

Keďže DNA a RNA molekuly sú veľmi malé, signál zaznamenaný počas nanopórového sekvenovania odráža interakcie jednotlivých nukleotidov (báz) s nanopórom \cite{Zheng2023Nanopore}.
Každá báza má jedinečný tvar a elektrické vlastnosti, ktoré ovplyvňujú spôsob, akým mení elektrický prúd pri prechode cez nanopór \cite{Zheng2023Nanopore}.
Preto je možné získať signál pre každú jednotlivú bázu v molekule, čo umožňuje presné čítanie sekvencie \cite{Zheng2023Nanopore}.
